\documentclass[a4paper,11pt]{article}
\usepackage[a4paper,margin=2cm]{geometry}
\usepackage[T1]{fontenc}
\usepackage[utf8]{inputenc}
\usepackage[ngerman,english]{babel}
\usepackage{lmodern}
\usepackage{microtype}
\usepackage{xcolor}
\usepackage{parskip}
\usepackage{enumitem}
\usepackage[most]{tcolorbox}
\usepackage{titlesec}
\usepackage[hidelinks]{hyperref}
\usepackage{array}
\usepackage{longtable}
\usepackage[table]{xcolor}
\definecolor{HeaderBlueDark}{RGB}{70,110,170}
\definecolor{RowBlueLight}{RGB}{235,243,255}
\definecolor{RowBlueDark}{RGB}{220,233,250}
\definecolor{SoftGray}{RGB}{90,90,90}
\setlength{\parindent}{0pt}
\setlist[itemize]{leftmargin=1.4em,itemsep=0.35em,topsep=0.35em}
\linespread{1.06}
\renewcommand{\arraystretch}{1.35}
\setlength{\tabcolsep}{5pt}
\titleformat{\section}[block]{\Large\bfseries\color{HeaderBlueDark}}{}{0pt}{}
\titlespacing{\section}{0pt}{1.3em}{0.8em}
\titleformat{\subsection}[block]{\large\bfseries\color{HeaderBlueDark}}{}{0pt}{}
\titlespacing{\subsection}{0pt}{1.0em}{0.6em}
\newtcolorbox{exbox}{enhanced,breakable,colback=RowBlueLight,colframe=RowBlueDark,boxrule=0.4pt,arc=1.5mm,left=1.2mm,right=1.2mm,top=1mm,bottom=1mm,title={Beispiele \textnormal{(Examples)}},fonttitle=\bfseries,colbacktitle=RowBlueDark,coltitle=black}
\NewDocumentEnvironment{grammarentry}{mm+b}{%
\begin{tcolorbox}[enhanced,breakable,colback=white,colframe=HeaderBlueDark,boxrule=0.6pt,arc=1.5mm,left=1.5mm,right=1.5mm,top=1mm,bottom=1mm,colbacktitle=HeaderBlueDark,coltitle=white,fonttitle=\bfseries,title={#1\hfill\normalfont\footnotesize #2}]
#3
\end{tcolorbox}}{}
\newcommand{\GermanText}[1]{\textbf{Deutsch: }#1\par}
\newcommand{\EnglishText}[1]{{\small\color{SoftGray}\textbf{English: }#1\par}}
\newcommand{\ExampleItem}[2]{\item \textbf{#1}\\{\small\color{SoftGray}#2}}
\newcommand{\DEEN}[2]{\textbf{#1}\par{\footnotesize\color{SoftGray}#2}}
\title{Vorgangspassiv\\[0.3em]\large\textnormal{Event passive}}
\date{}
\begin{document}
\maketitle
\tableofcontents
\newpage

\section{Grundidee -- Basic idea}
\begin{grammarentry}{Bedeutung}{Meaning}
\GermanText{Das Vorgangspassiv beschreibt eine \textbf{Handlung oder einen Vorgang}. Wichtig ist, was geschieht; der Handelnde ist oft unwichtig oder unbekannt.}
\EnglishText{The event passive describes an \textbf{action or process}. What happens is important; the agent is often unimportant or unknown.}
\GermanText{Bildung: \textbf{werden + Partizip II}.}
\EnglishText{Formation: \textbf{werden + past participle}.}
\end{grammarentry}
\begin{exbox}\begin{itemize}
\ExampleItem{Aktiv: Der Mechaniker repariert das Auto.}{Active: The mechanic repairs the car.}
\ExampleItem{Passiv: Das Auto wird repariert.}{Passive: The car is being repaired.}
\end{itemize}\end{exbox}

\section{Zeitformen -- Tenses}
\rowcolors{2}{RowBlueLight}{RowBlueDark}
\begin{longtable}{>{\bfseries}p{3.0cm}|p{5.2cm}|p{5.6cm}}
\rowcolor{HeaderBlueDark}\color{white}\textbf{Zeit / Tense}&\color{white}\textbf{Bildung / Formation}&\color{white}\textbf{Beispiel / Example}\\
\DEEN{Präsens}{Present}&\DEEN{werden + Partizip II}{werden + past participle}&\DEEN{Die Tür wird geöffnet.}{The door is being opened.}\\
\DEEN{Präteritum}{Simple past}&\DEEN{wurde + Partizip II}{wurde + past participle}&\DEEN{Die Tür wurde geöffnet.}{The door was opened.}\\
\DEEN{Perfekt}{Present perfect}&\DEEN{sein + Partizip II + worden}{sein + past participle + worden}&\DEEN{Die Tür ist geöffnet worden.}{The door has been opened.}\\
\DEEN{Plusquamperfekt}{Past perfect}&\DEEN{war + Partizip II + worden}{war + past participle + worden}&\DEEN{Die Tür war geöffnet worden.}{The door had been opened.}\\
\DEEN{Futur I}{Future I}&\DEEN{werden + Partizip II + werden}{werden + past participle + werden}&\DEEN{Die Tür wird geöffnet werden.}{The door will be opened.}\\
\DEEN{Futur II}{Future perfect}&\DEEN{werden + Partizip II + worden sein}{werden + past participle + worden sein}&\DEEN{Die Tür wird geöffnet worden sein.}{The door will have been opened.}\\
\end{longtable}

\section{Aktiv zu Passiv -- Active to passive}
\begin{grammarentry}{Akkusativobjekt wird Subjekt}{Accusative object becomes subject}
\GermanText{Das Akkusativobjekt des Aktivsatzes wird normalerweise zum Nominativsubjekt des Passivsatzes.}
\EnglishText{The accusative object of the active sentence normally becomes the nominative subject of the passive sentence.}
\GermanText{Aktiv: \textbf{Der Arzt untersucht den Patienten.} Passiv: \textbf{Der Patient wird vom Arzt untersucht.}}
\EnglishText{Active: \textbf{The doctor examines the patient.} Passive: \textbf{The patient is examined by the doctor.}}
\end{grammarentry}
\begin{grammarentry}{Dativ bleibt Dativ}{Dative remains dative}
\GermanText{Ein Dativobjekt bleibt beim Wechsel ins Passiv Dativ.}
\EnglishText{A dative object remains dative when changing to the passive.}
\GermanText{Aktiv: \textbf{Die Lehrerin erklärt dem Schüler die Regel.} Passiv: \textbf{Die Regel wird dem Schüler erklärt.}}
\EnglishText{Active: \textbf{The teacher explains the rule to the student.} Passive: \textbf{The rule is explained to the student.}}
\end{grammarentry}

\section{von und durch -- von and durch}
\begin{grammarentry}{von + Dativ}{Agent}
\GermanText{\textbf{von + Dativ} nennt häufig die handelnde Person, Institution oder unmittelbare Ursache.}
\EnglishText{\textbf{von + dative} often names the acting person, institution, or direct cause.}
\end{grammarentry}
\begin{grammarentry}{durch + Akkusativ}{Means or mediating cause}
\GermanText{\textbf{durch + Akkusativ} nennt häufig ein Mittel, einen Vorgang oder eine vermittelnde Ursache.}
\EnglishText{\textbf{durch + accusative} often names a means, process, or mediating cause.}
\end{grammarentry}
\begin{exbox}\begin{itemize}
\ExampleItem{Das Gesetz wurde vom Parlament beschlossen.}{The law was passed by parliament.}
\ExampleItem{Die Straße wurde durch den Sturm beschädigt.}{The road was damaged by the storm.}
\end{itemize}\end{exbox}

\section{worden und geworden -- worden and geworden}
\begin{grammarentry}{Wichtiger Unterschied}{Important distinction}
\GermanText{Im Perfekt und Plusquamperfekt des Vorgangspassivs steht \textbf{worden}, nicht \textbf{geworden}: \textbf{Das Haus ist gebaut worden.}}
\EnglishText{In the perfect and past perfect of the event passive, use \textbf{worden}, not \textbf{geworden}: \textbf{The house has been built.}}
\GermanText{\textbf{geworden} gehört zum Vollverb \textbf{werden} im Sinn von \glqq become\grqq: \textbf{Er ist Arzt geworden.}}
\EnglishText{\textbf{geworden} belongs to the full verb \textbf{werden} meaning ``become'': \textbf{He became a doctor.}}
\end{grammarentry}

\section{Merksatz -- Key rule}
\begin{grammarentry}{Kurz}{In short}
\GermanText{\textbf{werden + Partizip II} = Handlung/Vorgang im Passiv.}
\EnglishText{\textbf{werden + past participle} = action/process in the passive.}
\end{grammarentry}
\end{document}
